\begin{abstract}
Masticatory motor output depends on continuous periodontal and mucosal
somatosensory feedback mediated by trigeminal afference.
Whether the motor system compensates for transient reduction of this feedback
has not been rigorously characterised.
We present a secondary electromyographic (EMG) analysis of a within-subject
crossover dataset ($N$\,=\,34) in which unilateral topical lidocaine (10\%) or
matched placebo spray was applied before four 60-second chewing bouts occurring
at approximately 3, 7.5, 12, and 17~min post-application.
Masseter activity was recorded bipolarly at 1024~Hz, filtered 20--450~Hz with
Zapline-plus line-noise suppression, and analysed bout-by-bout using paired
Wilcoxon tests with Benjamini--Hochberg correction and repeated-measures ANOVA.
Session-averaged comparisons across all four bouts showed no significant
differences for any metric (equivalence confirmed by two one-sided tests at
$d$\,=\,0.5).
Resolving by bout revealed a time-dependent effect: at bout~1
($\approx$3~min post-application, maximal anesthesia), median cycle amplitude
was greater under anesthesia ($d_z$\,=\,0.695, $p$\,=\,0.0003,
Benjamini--Hochberg FDR-corrected $p$\,=\,0.0035), alongside increases in
RMS amplitude, total power, and spectral power in 20--150~Hz bands
(all FDR-corrected $p$\,<\,0.03) and greater inter-cycle regularity
(CV of inter-peak interval $d_z$\,=\,$-$0.434, FDR $p$\,=\,0.042).
These effects decayed monotonically across bouts~2--4 (trend $p$\,=\,0.0002),
consistent with topical lidocaine washout.
Critically, median frequency, chewing rate, and duty cycle showed no change at
any bout (non-significant at all time points; two one-sided test equivalent).
Cycle-locked time-frequency analysis identified a broadband cluster
($\approx$30--300~Hz) of greater EMG power under anesthesia at bout~1
(cluster permutation $p$\,=\,0.035 directional; two-tailed $p$\,=\,0.051),
absent at bout~4.
These results indicate that transient oral somatosensory deprivation elicits
sensorimotor gain compensation in the masseter -- greater amplitude and
regularity without timing change -- consistent with a hierarchical model in
which the brainstem central pattern generator maintains rhythm while peripheral
feedback modulates efferent gain.
This preserved rhythmic output provides a post-hoc manipulation check for the
parent study and a mechanistic account of why the cognitive and cortical effects
of chewing are unaffected by dental anesthesia.
\end{abstract}

% ─── BILINGUAL ABSTRACT (Spanish — for departmental records; not for SR submission) ───
%
% \begin{abstract}[Resumen]
% La actividad motora masticatoria depende de la retroalimentación somatosensorial
% periodontal y mucosa mediada por la vía trigeminal. No se ha caracterizado
% rigurosamente si el sistema motor compensa la reducción transitoria de esta
% retroalimentación. Presentamos un análisis electromiográfico secundario de un
% conjunto de datos cruzados intrasujeto ($N$\,=\,34) en el que se aplicó
% lidocaína tópica unilateral (10\%) o spray placebo equivalente antes de cuatro
% series de masticación de 60~s a $\approx$3, 7,5, 12 y 17~min post-aplicación.
% El análisis muestra que la amplitud del ciclo masticatorio fue mayor bajo
% anestesia en la primera serie ($d_z$\,=\,0{,}695, $p_{\text{FDR}}$\,=\,0{,}0035),
% sin cambios en frecuencia, tasa o ciclo de trabajo, con decaimiento consistente
% con el lavado farmacológico de la lidocaína. Estos resultados indican que la
% privación somatosensorial oral transitoria elicita una compensación del ganancia
% motora maseterina sin alterar el ritmo, consistente con un modelo jerárquico en
% el que el generador central de patrones del tronco encefálico mantiene el ritmo
% mientras la retroalimentación periférica modula la ganancia eferente.
% \end{abstract}

